% Verbatim from minor_report/chapters/chapter_3_methodology.tex lines 1655-1772 (retired report),
% headings demoted by 1 level(s). Do not edit here; edit the source of the numbers.

Detector precision, recall, TP, FP, FN, moving/static recall, IoU, and image-area
coverage are computed by \repo{yolo_eval.py} at recorded confidence, matching-IoU,
occlusion, minimum-size, and motion thresholds. Mask-use statistics are joined to
trajectories by \repo{timestamp_ns}. A semantic class is never treated as
ground-truth motion.

\paragraph{Constructing the ground-truth boxes} 
\label{sec:yolo-gt-projection}

The detector is scored against boxes computed from the simulator's own state
rather than against hand annotation. No image is labelled by hand anywhere in
this project. The chain from a recorded object pose to a box in pixels is as
follows; all of it is in \repo{yolo_eval.py} and \repo{yolo_gt.py}.

\begin{enumerate}[leftmargin=1.6em]
  \item \textbf{Geometry, not just pose.} A pose is a point and a box needs
    extent, so \repo{yolo_gt.vertices()} reads the prop's collision mesh
    (\repo{.dae}) from the Gazebo model directory and returns its vertices
    $V$ in the model frame. Meshes above \num{4000} vertices are decimated with
    a fixed seed: a few thousand points bound a silhouette as tightly as
    \num{100000} and keep the per-frame cost flat.
  \item \textbf{Object to world.} The pose $(R_{wo}, p_{wo})$ recorded for that
    object places the vertices: $v_w = R_{wo} V + p_{wo}$.
  \item \textbf{Camera pose.} The robot pose gives $T_{wb}$, and the camera is
    placed on the body by the extrinsic $T_{bc}$ read from
    \repo{config/wil_sim/stereo_imu.yaml}, so $T_{wc} = T_{wb} T_{bc}$. Reading
    the extrinsic from the estimator's own configuration file is deliberate: the
    evaluation and VINS-Fusion then cannot disagree about where the camera is.
  \item \textbf{World to camera}, by inverting that transform, using the
    transpose because the rotation is orthonormal:
    $R_{cw} = R_{wc}^{\top}$, $t_{cw} = -R_{cw} t_{wc}$, $v_c = R_{cw} v_w +
    t_{cw}$.
  \item \textbf{Projection.} Vertices with $z \le \SI{0.05}{\metre}$ are
    discarded \emph{before} the division, then the remainder are projected
    through an ideal pinhole:
    \begin{equation}
      u = f_x \frac{x_c}{z_c} + c_x, \qquad
      v = f_y \frac{y_c}{z_c} + c_y .
      \label{eq:pinhole}
    \end{equation}
  \item \textbf{Box.} The axis-aligned extent of the projected points,
    $[\min u, \max u] \times [\min v, \max v]$, clipped to the image.
\end{enumerate}

The $z$ guard in step 5 is not a formality. A vertex behind the camera has
negative $z$, and the division then mirrors it to the wrong side of the image,
producing a well-formed box in entirely the wrong place --- a failure that looks
like a detector error rather than a projection error. Objects that straddle the
image plane are projected from their visible vertices only.

Each object also stores $\tilde{z} = \mathrm{median}(z_c)$, which is what the
occlusion test below uses to decide which of two overlapping objects is nearer.

\begin{table}[htbp]
\centering
\caption[Parameters of the detection evaluation.]{Parameters of the detection evaluation. The intrinsics are not
calibrated: Gazebo renders an ideal pinhole, so $f_x = f_y = (W/2)/\tan(\phi/2)$
with every distortion term zero, and the values below are exact rather than
estimated. The three exemption rules decide which ground-truth objects are scored
at all, and are the reason the recall in \cref{sec:yolo-performance} is computed
over a subset of the objects present.}
\label{tab:yolo-eval-params}
\small
\begin{tabularx}{\textwidth}{>{\raggedright\arraybackslash}p{0.30\textwidth}lX}
\toprule
Parameter & Value & Role \\
\midrule
$f_x$, $f_y$ & \num{640}, \num{640} & ideal pinhole, \SI{90}{\degree} horizontal
  field of view \\
$c_x$, $c_y$ & \num{640}, \num{360} & principal point at image centre \\
Image size & $1280 \times 720$ & \\
Distortion & all zero & the renderer applies none \\
Mesh vertices per prop & $\le \num{4000}$ & decimated with a fixed seed \\
Near-plane guard & \SI{0.05}{\metre} & vertices nearer than this are dropped
  before projection \\
\midrule
Confidence, \repo{--conf} & \num{0.35} & detections below this are never emitted \\
NMS IoU, \repo{--iou-nms} & \num{0.5} & \\
Matching IoU, \repo{--iou-match} & \num{0.5} & prediction counts as a hit above
  this \\
Motion threshold, \repo{--move-thresh} & \SI{0.05}{\metre\per\second} & object
  is ``moving'' above this, from its pose history \\
\midrule
\textbf{Exemption:} minimum size, \repo{--min-box-px} & \SI{12}{px} & smaller
  ground-truth boxes are discarded and never appear \\
\textbf{Exemption:} occlusion, \repo{--occlusion} & \num{0.70} & ground truth
  covered more than this by nearer objects is scored as neither hit nor miss
  (\emph{yellow} in \cref{fig:yolo-ground-truth-overlay}) \\
\textbf{Exemption:} scenery cover, \repo{--ignore-cover} & \num{0.50} & a
  prediction this covered by unlabelled scenery is neither credited nor penalised \\
\bottomrule
\end{tabularx}
\end{table}

\paragraph{What this convention costs} 

Two consequences follow from scoring against projected geometry, and both bear on
how the detector numbers in \cref{sec:yolo-performance} should be read.

The boxes are \emph{tighter than a human would draw}. They are the exact
silhouette of the collision mesh, whereas an annotator draws a looser rectangle
and rounds generously. A detector trained on human boxes is therefore scored
against a stricter convention than it was fitted to, and loses IoU that
represents no error on its part. Of the differences between how this detector
is scored here and how it was scored during training, this is the one that cannot
be removed without relabelling.

The exemptions are \emph{not symmetric with the training convention}. A
human-annotated set has no notion of an exempt object: every labelled box counts.
The occlusion and minimum-size rules exist because a crate that is nine tenths
behind a shelf is not a fair test of a detector, which is the right decision for
reporting detector quality --- but it means recall here is computed over the
easier subset. The size of that effect is measurable from the stored counts:
\num{34476} of the \num{96758} ground-truth objects across the five recordings
are exempt, and recomputing recall with every one of them counted as a miss moves
it from \num{0.593} to \num{0.382}
(\repo{script/yolo_generalization.py}). Any comparison against a
training-set recall has to use the second figure, not the first.
